\section{引言}
\label{sec:intro}

逐次仅开启一个光源拍摄物体（OLAT），再从新视角对其重光照，是构建照片级真实数字资产的经典途径之一~\cite{debevec2000acquiring}。
如今，这一流程通常从数百至数千张图像重建可重光照表示，每张图像由单个经过标定的点光源照明~\cite{bi2020neural,zeng2023nrhints}；高斯泼溅~\cite{kerbl20233d} 则使重建结果能够实时渲染~\cite{bi2024gs3,kuang2024olat,fan2025rng,liu2025bigs,zheng2026ssdgs}。

在点光源照明下，图像是否真实，很大程度上在材质求值之前就已决定：光照到哪里，又被挡在了哪里。
现有高斯方法逐个图元处理这一问题。
\gsthree~\cite{bi2024gs3} 和 SSD-GS~\cite{zheng2026ssdgs} 从光源视角泼溅场景，为每个高斯赋予一个可见性值；RNG~\cite{fan2025rng} 将深度阴影图归结为一个标量提示；半透明效果则由逐高斯网络根据局部输入预测~\cite{dihlmann2024sss,zheng2026ssdgs}。
因此，地面上的一个大型高斯只能承载单一阴影值，移动阴影的每一道边界都必须依靠增密来刻画。

我们转而在光源的成像平面上进行建模。
从点光源视角光栅化高斯时，alpha 合成权重将每个光源像素的通量分配给沿其射线排列的图元；在泼溅模型内，这些权重就是各高斯截获该通量的比例。
因此，只需额外进行一次光栅化，即可记录每个图元吸收了多少光、吸收发生在何种深度；我们将这一记录称为\emph{沉积量}（deposit）。
正向渲染早已通过阴影图~\cite{williams1978casting,donnelly2006variance,peters2015moment}、半透明阴影图~\cite{dachsbacher2003translucent} 和反射阴影图~\cite{dachsbacher2005reflective} 利用光源视角，在已知几何上存储人工设计的量。
我们将这一思路引入逆向渲染，同时学习几何及图中存储的量。

我们的方法 \ours（Light-Space Atlas，光源空间图集）针对每个光源，将沉积量的前两个深度矩和少量可学习通量特征泼溅到图集中，再经滤波构成金字塔（\cref{fig:teaser}）。
任何需要进行着色求值的点都称为\emph{接收点}：它投影到光源视图并读取金字塔；几何优化期间使用高斯中心，之后使用像素对应的表面点。
深度矩提供一种类似方差阴影图~\cite{donnelly2006variance} 的滤波阴影检验，再由残差网络修正；通量特征则驱动一个对沉积特征呈线性的光传输项，可视为半透明阴影图的可学习变体。
图集为每张图像增加一次 $512^2$ 光栅化；其输入均相对于光源定义，且梯度能够通过图集传回高斯。

若使用显示编码辐亮度上的损失训练该模型，Lego 推土机铲斗等细小、密集结构会退化为团块（\cref{fig:loss_domain}）。
为此，我们按辐射度量域分阶段优化：首先用不含网络的点光源模型拟合几何，同时将目标从显示编码值逐步退火到线性辐亮度，类似于 \gsthree 和 SSD-GS 对 HDR 数据的处理~\cite{bi2024gs3,zheng2026ssdgs}；随后，以高斯为接收点在线性辐亮度域训练完整模型；最后冻结几何，以像素为接收点精修外观。
在两个场景的受控实验中，仅改变损失计算域，就使留出视图的 PSNR 从 Lego 上的 23.0 提升至 30.2\,dB，从 Drums 上的 24.1 提升至 30.1\,dB。

我们在三个公开 OLAT 基准~\cite{zeng2023nrhints,bi2024gs3,dihlmann2024sss} 的全部 18 个场景上，按统一协议重新训练 \gsthree、SSD-GS、OLAT Gaussians~\cite{kuang2024olat} 和 RNG，并对照完全相同的目标图像评分。
\ours 获得了最佳平均 PSNR、SSIM 和 LPIPS（PSNR 为 33.47\,dB，次优方法为 31.59\,dB），每个场景训练约需 20 分钟。
分基准来看，它在真实拍摄数据上领先最强基线 2.0\,dB（逐视图二维对齐后为 0.7\,dB），在 \gsthree 场景上领先 0.4\,dB，在 SSS-GS 上落后 OLAT Gaussians 0.1\,dB，因而成为唯一在每个基准上都位列前二的方法。
消融实验也揭示了光传输项的局限：它在半透明材质上作用最大，而在多数其他材质上，相同规模的局部网络即可达到相当效果。
本文的贡献包括：
\begin{itemize}
\item 提出面向高斯泼溅的点光源可见性与光传输的光源空间表述。一次光源视角光栅化即可记录沉积量，接收点据此获得经滤波、可学习的阴影检验，以及对沉积特征呈线性的光传输项；
\item 提出按辐射度量域分阶段的优化流程，使该模型能够学习细小、密集的几何结构；受控实验表明，损失计算域决定了这些几何结构能否形成；
\item 在 18 个场景上分析光源空间传输何时有益、何时局部网络已经足够，并指出现有 OLAT 基准中测试光源与训练光源的最近夹角中位数仅为 $1.1^\circ$，因此这些基准几乎没有检验光源外推能力。
\end{itemize}
